%==============================================================================
%  CHAPTER 7 — SPRINT 4: OBSERVABILITY & PUBLIC SERVING
%==============================================================================
\chapter{Sprint 4 --- Observability and Public Serving}
\label{chap:sprint4}

\section*{Introduction}
The final sprint makes the platform observable and opens deployed models to the
outside world. It adds a live metrics dashboard, per-deployment public \gls{api}
keys with an unauthenticated prediction endpoint, and a set of diagnostics that
make failed runs easy to understand.

% ---------------------------------------------------------------------------
\section{Sprint Backlog}
\begin{table}[H]\centering
\caption{Sprint~4 backlog}\label{tab:backlog-s4}
\begin{tabularx}{\textwidth}{p{0.6cm} X p{2cm}}
\toprule
\textbf{\#} & \textbf{Task} & \textbf{Pts}\\
\midrule
12 & Collect cluster and per-pod resource metrics & 3\\
12 & Build the monitoring dashboard & 2\\
13 & Generate and store per-deployment API keys (hashed) & 3\\
13 & Expose a public, key-authenticated prediction endpoint & 2\\
14 & Capture and surface user-script errors for failed runs & 3\\
14 & Add a static pre-flight code analyzer with warnings & 2\\
\bottomrule
\end{tabularx}
\end{table}

% ---------------------------------------------------------------------------
\section{Analysis and Conception}

\subsection{Use Cases}
This sprint covers \emph{View metrics dashboard}, \emph{Create / revoke API keys},
\emph{Call the public prediction endpoint} (external client), and
\emph{Diagnose a failed run}.

\subsection{Sequence Diagram: Public Prediction with an API Key}
% >>> USER: render diagrams/seq_public_predict.puml -> images/seq_public_predict.pdf
\begin{figure}[H]\centering
  % \includegraphics[width=\textwidth]{seq_public_predict.pdf}
  \fbox{\parbox[c][5.5cm][c]{\textwidth}{\centering\itshape
    Placeholder --- insert \texttt{images/seq\_public\_predict.pdf}\\
    (render \texttt{diagrams/seq\_public\_predict.puml})}}
  \caption{Sequence diagram of an external client calling the public endpoint}
  \label{fig:seq-public}
\end{figure}

\subsection{Data Model Increment}
The sprint adds the \texttt{DeploymentApiKey} entity (hashed key, display prefix,
name, timestamps, revocation marker) linked to a deployment.

% ---------------------------------------------------------------------------
\section{Realization}

\subsection{Observability Dashboard}
The backend queries the cluster's metrics service for per-pod \gls{cpu} and memory
usage and aggregates cluster-wide figures. The dashboard surfaces these as live
gauges and charts so the user can see the health of the cluster and of each
deployed model.

\subsection{Public Prediction API with API Keys}
Deploying a model (Sprint~3) makes it reachable from \emph{inside} the platform.
Sprint~4 takes the final step: it lets the owner expose that same model to the
\emph{outside world} so it can be consumed by a third-party website, a mobile
application, or any external backend, \textbf{without} that consumer needing a
platform account. This turns a trained model into a genuinely reusable web
service.

The mechanism is \textbf{per-deployment \gls{api} keys}. Once a deployment is in
the \texttt{READY} state, the owner opens it and clicks \emph{New key}; the
platform generates a key of the form \texttt{mlops\_<random>} and \textbf{shows
it once}. Only a \gls{sha}-256 hash of the key is stored, together with a short
display prefix (e.g.\ \texttt{mlops\_SdHvJo\dots}) so the owner can recognize it
later; the plaintext is never persisted and never shown again. A key can be
\emph{revoked} at any time, which immediately invalidates it.

External clients then call a dedicated \textbf{public endpoint},
\texttt{POST /api/public/predict/\{deployment\_id\}}, which lives \emph{outside}
the \gls{jwt}-protected part of the \gls{api}. The request carries the key in a
standard \texttt{Authorization: Bearer mlops\_<key>} header. The backend hashes
the presented key, looks it up scoped to that deployment, confirms the deployment
is ready, and forwards the request to the model server, returning the prediction.

\begin{lstlisting}[language=Python,caption={The public endpoint authenticates by hashed API key before forwarding.},label={lst:public-predict}]
key_hash = sha256(raw_key.encode()).hexdigest()
api_key = await lookup_active_key(db, key_hash, deployment_id)
if api_key is None:
    raise HTTPException(401, "Invalid or revoked API key")
return deployment_service.predict(deployment.inference_service_name, payload.instances)
\end{lstlisting}

To make integration effortless, the deployment page includes a \textbf{code-snippet
generator}: once a key exists, it shows ready-to-use \textbf{cURL}, \textbf{JavaScript
(fetch)}, and \textbf{Python (requests)} snippets, pre-filled with the deployment's
public URL and the \texttt{Authorization} header. The user copies one snippet,
substitutes the key they saved at creation, and pastes it directly into their own
website or service, as shown in Listing~\ref{lst:public-curl}.

\begin{lstlisting}[language=bash,caption={Auto-generated cURL snippet a user can paste into any website or service.},label={lst:public-curl}]
curl -X POST "https://mlops.firasmahjoubi.app/api/public/predict/<deployment_id>" \
  -H "Authorization: Bearer mlops_<your_key>" \
  -H "Content-Type: application/json" \
  -d '{"instances": [[5.1, 3.5, 1.4, 0.2]]}'
\end{lstlisting}

\subsection{Administration}
The platform distinguishes two roles through a \texttt{role} field on the user
account (\texttt{user} or \texttt{admin}); the first administrators are
bootstrapped from an \texttt{ADMIN\_EMAILS} configuration list. A dedicated
backend dependency, \texttt{require\_admin}, guards every administrative endpoint
and returns \texttt{403} for non-administrators, while the Angular routes are
protected by an \texttt{adminGuard} and the administration menu is shown only to
administrators.

The administration area provides three capabilities:
\begin{itemize}[leftmargin=1.4em]
  \item \textbf{User management}, list all accounts and activate, deactivate,
        promote, demote, purge, or delete them (with safeguards that prevent an
        administrator from locking out their own account);
  \item \textbf{Platform overview}, aggregate counts of users, projects, pipeline
        runs, models, and deployments, together with cluster health;
  \item \textbf{Cross-user oversight}, read-only listings of every user's
        projects, deployments, and pipeline runs, each annotated with its owner.
\end{itemize}

% >>> USER: screenshot of the admin Users / Overview pages -> images/
\subsection{Diagnostics and Robustness}
Because uploaded code frequently fails for environmental reasons, the platform
captures the user-script error of a failed run and stores it durably, so the
interface can show a concise ``why did this fail'' explanation independent of the
underlying job logs. A static pre-flight analyzer also inspects uploaded code and
warns about common pitfalls, missing training calls, hard-coded local paths,
command-line argument usage, and notebook-only constructs, before the run is
launched.

\subsection{User Interface}
% >>> USER: screenshots: dashboard, API keys panel, "Why?" error panel.
The interface adds a metrics dashboard, an \gls{api}-keys management panel with a
code-snippet generator, and a failure-explanation panel on the runs list.

% ---------------------------------------------------------------------------
\section{Sprint Review and Retrospective}
At the end of Sprint~4 the platform is observable, its models are callable by
external clients, and failures are diagnosable, completing the end-to-end
\gls{mlops} loop.
\begin{figure}[H]\centering
\begin{tikzpicture}
\begin{axis}[
    width=0.8\textwidth, height=6cm,
    xlabel={Day}, ylabel={Remaining points},
    xmin=0, xmax=11, ymin=0, ymax=16,
    legend pos=north east, grid=both,
]
\addplot[dashed,gray] coordinates {(0,15)(10,0)};
\addplot[thick,primary,mark=*] coordinates {
  (0,15)(1,15)(2,13)(3,12)(4,10)(5,8)(6,6)(7,5)(8,3)(9,1)(10,0)};
\legend{Ideal, Actual}
\end{axis}
\end{tikzpicture}
\caption{Sprint~4 burndown chart}\label{fig:burndown-s4}
\end{figure}

\noindent\textbf{Retrospective.} % >>> USER
Diagnosing real-world dependency conflicts in user code proved a recurring theme;
investing in clear error capture paid off. \todo{your note}.
